\section{Introduction}
Effective research requires finding the right papers, understanding them
quickly, organizing them reliably, and collaborating without friction. In
practice, researchers juggle four to five disconnected tools --- reference
managers, cloud drives, PDF readers, note-taking apps, and general chatbots
--- losing context at every switch. Cognitive-load research and the
ScholarFlow survey agree on the cost: the average researcher uses 4.2
disconnected tools daily, context-switching overhead cuts productivity by
30--50\%, and current tool satisfaction sits at only 3.28 out of 5.

ScholarFlow was developed to consolidate this fragmented workflow into one
AI-native hub: upload and organize papers, read and annotate them, write with
integrated citations, and collaborate in shared workspaces with real-time
co-editing. Unlike legacy tools that treat AI as an optional add-on,
ScholarFlow makes AI a first-class capability of every workflow --- semantic
search instead of keyword search, generated summaries and Q\&A instead of
manual skimming, and contextual writing assistance instead of a detached
chatbot window.

This document serves four audiences. The \emph{development team} finds
architectural blueprints and implementation rules for the multi-tier system
(Next.js frontend, Express backend, PostgreSQL with pgvector). \emph{Project
evaluators} find scope boundaries, phased delivery evidence, and measurable
acceptance criteria. The \emph{quality assurance} view is covered by testable
requirements and performance budgets (sub-300\,ms semantic search, $<$50\,ms
per server-side query event, 99.9\% availability target). Finally,
\emph{stakeholders} find market validation --- a 32-respondent survey with
75\% need confirmation and 100\% AI-feature interest --- together with a
freemium revenue model designed for student affordability.

Development followed an Agile-iterative life cycle across phased rollouts:
an MVP slice (authentication $\rightarrow$ upload $\rightarrow$ library
$\rightarrow$ editor $\rightarrow$ sharing) first, then AI capabilities on
top of the library, then collaboration and real-time features, then
monetization and administration, closing with hardening and launch readiness.
Each phase shipped a thin vertical slice rather than broad incomplete
scaffolding, with continuous gates --- build, type-check, lint, and tests ---
required green on every merge. The project applies core software engineering
principles throughout: modularity (decoupled frontend/backend monorepo),
scalability (pooled connections, paginated and access-scoped queries),
maintainability (strict TypeScript, Zod contracts, migration discipline), and
security (JWT sessions, RBAC, verified webhooks) --- making it suitable for
real-world pilot deployment and iterative growth.
